% Partie : logique intuitionniste
% Chapitre : propositions comme types
% Section : réduction

\documentclass[../../../include/open-logic-section]{subfiles}

\begin{document}

\olfileid[fr]{int}{pty}{red}

\olsection{Réduction}

Dans une dérivation en déduction naturelle, une règle
d'introduction immédiatement suivie de la règle
d'élimination correspondante forme un détour que l'on
peut « réduire ». Par exemple, une \Intro{\land}
suivie de $\Elim{\land}$ peut être supprimée :
la conclusion de l'élimination est l'un des conjoints
introduits, et la preuve de ce conjoint figure déjà
parmi les prémisses de l'introduction.
On peut donc utiliser directement cette sous-preuve :
\begin{gather*}
  \bottomAlignProof
  \AxiomC{}
  \RightLabel{$\delta_1$}
  \DeduceC{$!A_1$}
  \AxiomC{}
  \RightLabel{$\delta_2$}
  \DeduceC{$!A_2$}
  \RightLabel{$\Intro{\land}$}
  \BinaryInfC{$!A_1 \land !A_2$}
  \RightLabel{$\Elim{\land}_i$}
  \UnaryInfC{$!A_i$}
  \DisplayProof\bottomAlignProof
  \quad
  \redone
  \quad
  \AxiomC{}
  \RightLabel{$\delta_i$}
  \DeduceC{$!A_i$}
  \DisplayProof
\end{gather*}

Cette simplification suggère une relation de réduction
sur les termes de preuve correspondants. Si $N_1$
est le terme de preuve de~$\delta_1$ et $N_2$
celui de~$\delta_2$, le terme de la preuve de gauche
se réduit en une étape au terme de la preuve de droite :
\[
  \proj{i}{\pair{N_1}{N_2}} \redone N_i
\]
Dans le lambda-calcul typé avec produits, c'est la
règle de $\beta$-réduction du type produit.

Un tel couple introduction--élimination est ici
appelé une \emph{coupure}, ou un détour réductible.
Dans le lambda-calcul typé, le terme à gauche de
$\redone$ est un \emph{redex} ; celui à droite en est
le \emph{réduit}. Dans le $\lambd$-calcul non typé pur,
les redex $\beta$ sont les applications d'une
abstraction à un argument, comme
$(\lambd[x][N])Q$, où la variable $x$ est liée dans $N$.
Notre syntaxe typée comporte
en outre $\pair{N}{M}$, $\proj{i}{N}$,
$\inj[!A]{i}{N}$, $\dcase{N}{x_1}{M_1}{x_2}{M_2}$
et $\abort{!A}{N}$ ; elle possède les redex
correspondant aux constructeurs et destructeurs
ajoutés. Le type vide n'ajoute toutefois aucun
redex introduction--élimination, faute de constructeur.

La disjonction possède une réduction analogue :
\begin{gather*}
  \bottomAlignProof
  \AxiomC{}
  \RightLabel{$\delta$}
  \DeduceC{$!A_i$}
  \RightLabel{$\Intro{\lor}$}
  \UnaryInfC{$!A_1 \lor !A_2$}
  \AxiomC{$\Discharge{!A_1}{x_1}$}
  \RightLabel{$\delta_1$}
  \DeduceC{$!C$}
  \AxiomC{$\Discharge{!A_2}{x_2}$}
  \RightLabel{$\delta_2$}
  \DeduceC{$!C$}
  \DischargeRule{$\Elim{\lor}$}{x_1,x_2}
  \TrinaryInfC{$!C$}
  \DisplayProof\bottomAlignProof
  \quad
  \redone
  \quad
  \AxiomC{}
  \RightLabel{$\delta$}
  \DeduceC{$!A_i$}
  \RightLabel{$\delta_i$}
  \DeduceC{$!C$}
  \DisplayProof
\end{gather*}
Elle correspond à la réduction suivante des termes de preuve :
\begin{gather*}
  \dcase{\inj[!A_{3-i}]{i}{M}}{x_1}{N_1}{x_2}{N_2} \redone \Subst{N_i}{M}{x_i}
\end{gather*}
où $M$ est le terme de preuve de~$\delta$ et $N_i$
celui de~$\delta_i$. C'est la règle de $\beta$-réduction
du type somme. Dans $\Subst{N_i}{M}{x_i}$,
on remplace les occurrences de $x_i$ libres dans
le terme de la branche~$N_i$ par~$M$,
en évitant toute capture de variables.

Pour le type fonction, la réduction supprime
le détour $\Intro\lif$ suivi de $\Elim\lif$ :
\begin{gather*}
  \bottomAlignProof
  \AxiomC{$\Discharge{!A}{x}$}
  \RightLabel{$\delta$}
  \DeduceC{$!B$}
  \DischargeRule{$\Intro{\lif}$}{x}
  \UnaryInfC{$!A \lif !B$}
  \AxiomC{}
  \RightLabel{$\delta'$}
  \DeduceC{$!A$}
  \RightLabel{$\Elim{\lif}$}
  \BinaryInfC{$!B$}
  \DisplayProof\bottomAlignProof
  \quad
  \redone
  \quad
  \AxiomC{}
  \RightLabel{$\delta'$}
  \DeduceC{$!A$}
  \RightLabel{$\delta$}
  \DeduceC{$!B$}
  \DisplayProof
\end{gather*}
Sur les termes de preuve, c'est la $\beta$-réduction usuelle :
\begin{gather*}
  (\lambd[\typeof{x}{!A}][N]M)
  \redone \Subst{N}{M}{x}
\end{gather*}

Outre ces conversions qui suppriment un détour,
on peut considérer des \emph{conversions de permutation}.
Elles s'appliquent, par exemple, lorsqu'une
\Elim{\lor} est suivie d'une autre règle d'élimination :
\begin{multline*}
  \bottomAlignProof
  \AxiomC{}
  \RightLabel{$\delta$}
  \DeduceC{$!A_j$}
  \RightLabel{$\Intro{\lor}[j]$}
  \UnaryInfC{$!A_1 \lor !A_2$}
  \AxiomC{$\Discharge{!A_1}{x_1}$}
  \RightLabel{$\delta_1$}
  \DeduceC{$!C_1 \land !C_2$}
  \AxiomC{$\Discharge{!A_2}{x_2}$}
  \RightLabel{$\delta_2$}
  \DeduceC{$!C_1 \land !C_2$}
  \DischargeRule{$\Elim{\lor}$}{x_1,x_2}
  \TrinaryInfC{$!C_1 \land !C_2$}
  \RightLabel{\Elim{\land}[i]}
  \UnaryInfC{$!C_i$}
  \DisplayProof\bottomAlignProof
  \quad
  \redone
  \\
  \AxiomC{}
  \RightLabel{$\delta$}
  \DeduceC{$!A_j$}
  \RightLabel{$\Intro{\lor}[j]$}
  \UnaryInfC{$!A_1 \lor !A_2$}
  \AxiomC{$\Discharge{!A_1}{x_1}$}
  \RightLabel{$\delta_1$}
  \DeduceC{$!C_1 \land !C_2$}
  \RightLabel{\Elim{\land}[i]}
  \UnaryInfC{$!C_i$}
  \AxiomC{$\Discharge{!A_2}{x_2}$}
  \RightLabel{$\delta_2$}
  \DeduceC{$!C_1 \land !C_2$}
  \RightLabel{\Elim{\land}[i]}
  \UnaryInfC{$!C_i$}
  \DischargeRule{$\Elim{\lor}$}{x_1,x_2}
  \TrinaryInfC{$!C_i$}
  \DisplayProof
\end{multline*}
Dans ce schéma, $M$ est le terme de la disjonction majeure
$!A_1 \lor !A_2$, introduction comprise. Les termes $N_1$ et $N_2$
sont ceux des branches. La conversion sur les termes lambda est :
\[
  \proj{i}{\dcase{M}{x_1}{N_1}{x_2}{N_2}} \redone 
\dcase{M}{x_1}{\proj{i}{N_1}}{x_2}{\proj{i}{N_2}}
\]

Les détours peuvent se trouver à l'intérieur d'une preuve :
on applique alors la conversion à la sous-preuve concernée.
De même, une réduction élémentaire s'applique à un sous-terme.
Si $N$ est un sous-terme de $M$ et $N \redone N'$,
remplacer cette occurrence de $N$ dans~$M$ par $N'$
produit un terme $M'$ tel que $M \redone M'$.
Nous écrivons $M \red M'$ lorsqu'il existe une suite
$M = M_1\redone M_2 \redone \dots \redone M_k = M'$ ;
le cas $k=1$ désigne la suite vide et donc $M=M'$.
Cette seconde relation est ainsi la clôture réflexive et
transitive de~$\redone$.

Une preuve à laquelle aucune conversion ne s'applique,
ni dans une sous-preuve, est dite \emph{normale}.
De même, un terme lambda dont aucun sous-terme
n'admet de conversion est \emph{normal}.
Un terme lambda normal est aussi une \emph{forme normale}.

\end{document}
